\usepackage{amsmath,amssymb,amsthm, amsfonts,mathrsfs}
\usepackage{geometry}
\usepackage{enumerate}
\usepackage{pdfpages}
\usepackage{subcaption}
\usepackage{graphicx}
\usepackage{cite}
\usepackage{tikz-cd}
\usepackage{tikz}
\usepackage{stix} %symbols and text fonts:formulas more clear
\usepackage{fancyhdr}
\usepackage{setspace}
\usepackage[many]{tcolorbox}
\usepackage[dvipsnames]{xcolor}                             % Link color
\usepackage{eso-pic} % 用于添加背景
\usepackage{mdframed}
% \usepackage{esint} %积分符号
\newcommand{\aint}{
	\mbox{
		\tikz[baseline = (int.base)] { \node (int) {$\int$}; \node at (int) {$-$}; }
	}
}

\newcommand\tbbint{{-\mkern -16mu\int}}
\newcommand\tbint{{\mathchar '26\mkern -14mu\int}}
\newcommand\dbbint{{-\mkern -19mu\int}}
\newcommand\dbint{{\mathchar '26\mkern -18mu\int}}
\newcommand\bint{
	{\mathchoice{\dbint}{\tbint}{\tbint}{\tbint}}
}
\newcommand\bbint{
	{\mathchoice{\dbbint}{\tbbint}{\tbbint}{\tbbint}}
}

\linespread{1.2} % 1.4-1.5x spacing

\usepackage{etoolbox} % for hooks
% Counter for numbered paragraphs
\newcounter{parnum}[section]
\renewcommand{\theparnum}{(\thesection.\arabic{parnum})}

% Define the numbered paragraph environment
\newcommand{\np}{\refstepcounter{parnum}\noindent\textbf{\theparnum}\quad}


\usepackage{charter}  % context less tails
\usepackage[T1]{fontenc}

\usepackage{xcolor}                             % Link color
 \definecolor{custom-blue}{RGB}{0,99,166} 
\usepackage{hyperref}
\hypersetup{
    colorlinks=true, % 启用彩色链接
    linkcolor=purple,  % 普通内部链接的颜色
    citecolor=blue,  % 引用链接的颜色
    filecolor=green,  % 打开本地文件的链接颜色
    urlcolor=red    % 链接到URL的颜色
}

\usepackage{titlesec}  % 用于修改章节标题的格式

% 使用 \titleformat 定义新的 \section 和 \subsection 样式
\titleformat{\section}[block]
  {\normalfont\Large\bfseries}    % 标题字体样式
  {\S\ \thesection}               % 在标题前加上 \S 和节编号
  {1em}                           % 标题和编号之间的间距
  {}                               % 标题后面的内容

\titleformat{\subsection}[block]
  {\normalfont\large\bfseries}    % 子节标题字体样式
  {\S\ \thesubsection}            % 在子节标题前加上 \S 和子节编号
  {1em}                           % 标题和编号之间的间距
  {}                               % 标题后面的内容

  \titleformat{\chapter}[block]
  {\centering\Huge\bfseries}    
  {Chapter \thechapter }           % 居中
  {1em}                           
  {} 

%%%%%%%%页面与页面对齐
\textwidth=6.5in
\textheight=8.7in
\topmargin=0in
\oddsidemargin=0in
\evensidemargin=0in



%%Bold the titles
\makeatletter
\def\th@plain{%
  \thm@notefont{}% same as heading font
  \itshape % body font
}
\def\th@definition{%
  \thm@notefont{}% same as heading font
  \normalfont % body font
}
\makeatother

\newtheorem{theorem}{Theorem}[section]
\newtheorem*{theorem*}{Theorem}
\newtheorem{lemma}[theorem]{Lemma}
\newtheorem{corollary}[theorem]{Corollary}
\newtheorem{proposition}[theorem]{Proposition}
\newtheorem{conjecture}[theorem]{Conjecture}

\theoremstyle{definition}
\newtheorem{definition}[theorem]{Definition}
\newtheorem{example}[theorem]{Example}
\newtheorem{exercise}[theorem]{Exercise}
\newtheorem{problem}[theorem]{Problem}
\newtheorem{remark}{Remark}[section]
\newtheorem{question}{Question}[section]
\renewcommand*{\thequestion}{\Alph{question}}

\DeclareMathOperator{\grad}{grad}
\DeclareMathOperator{\dimension}{dim}
\DeclareMathOperator{\trace}{tr} 
\DeclareMathOperator{\di}{div}
\DeclareMathOperator{\ricci}{Ric}
\DeclareMathOperator{\riem}{Riem}
\DeclareMathOperator{\dist}{dist}
\DeclareMathOperator{\hess}{Hess}
\DeclareMathOperator{\Id}{Id}
\DeclareMathOperator{\sing}{Sing}
\DeclareMathOperator{\ind}{Ind}

\DeclareMathOperator{\im}{Im}

\def\zz{{\mathbb Z}}
\def\rr{{\mathbb R}}
\def\cc{{\mathbb C}}
\def\qq{{\mathbb Q}}
\def\nn{{\mathbb N}}
\def\ss{{\mathbb S}}
\def\oo{{\infty}}
\def\aa{{\mathit A}}

\newcommand{\la}{\langle}
\newcommand{\ra}{\rangle}
\newcommand{\lap}{\Delta}


\setlength{\headheight}{14.5pt}




  \definecolor{headblue}{RGB}{0,102,204}  % 你可以换成你喜欢的颜色
  \pagestyle{fancy}
  \fancyhf{}
  \renewcommand{\sectionmark}[1]{\markboth{#1}{}}  % 更新 leftmark
  \lhead{\textcolor{headblue}{\nouppercase{\leftmark}}}  % 设置颜色
  \rhead{\thepage}
  
  
  \onehalfspacing
  
  \newcommand{\transparent}[1]{\color{black!#1}} % 定义替代命令
  \definecolor{bgcolor}{RGB}{248, 248, 255} % 添加这行定义背景色
  % 背景设置
  \AddToShipoutPictureBG{%
    \AtPageLowerLeft{%
      \color{bgcolor}%
      \rule{\paperwidth}{\paperheight}% 底色
    }%
    \AtPageCenter{%
      \put(0,-300){%
        \transparent{0.05}% 透明度
      }%
    }%
  }  

% 自定义定理框样式
\mdfdefinestyle{theoremstyle}{
  backgroundcolor=blue!5,
  linecolor=blue!40,
  linewidth=1pt,
  topline=false,
  bottomline=false,
  rightline=false,
  leftline=true,
  leftmargin=10pt,
  skipabove=10pt,
  skipbelow=10pt,
  innertopmargin=5pt,
  innerbottommargin=5pt,
  innerleftmargin=10pt,
  innerrightmargin=10pt
}


% 定义类：黄色
\mdfdefinestyle{definitionstyle}{
  backgroundcolor=yellow!3,  
  linecolor=yellow!40,
  linewidth=1pt,
  topline=false, bottomline=false, rightline=false, leftline=true,
  leftmargin=10pt,
  skipabove=10pt, skipbelow=10pt,
  innertopmargin=5pt, innerbottommargin=5pt, innerleftmargin=10pt, innerrightmargin=10pt
}

% 备注类：灰色
\mdfdefinestyle{remarkstyle}{
  backgroundcolor=gray!10,
  linecolor=gray!50,
  linewidth=1pt,
  topline=false, bottomline=false, rightline=false, leftline=true,
  leftmargin=10pt,
  skipabove=10pt, skipbelow=10pt,
  innertopmargin=5pt, innerbottommargin=5pt, innerleftmargin=10pt, innerrightmargin=10pt
}

% 练习类：橙色
\mdfdefinestyle{exercisestyle}{
  backgroundcolor=orange!10,
  linecolor=orange!50!black,
  linewidth=1pt,
  topline=false, bottomline=false, rightline=false, leftline=true,
  leftmargin=10pt,
  skipabove=10pt, skipbelow=10pt,
  innertopmargin=5pt, innerbottommargin=5pt, innerleftmargin=10pt, innerrightmargin=10pt
}

% 问题类：红色
\mdfdefinestyle{problemstyle}{
  backgroundcolor=red!5,
  linecolor=red!50,
  linewidth=1pt,
  topline=false, bottomline=false, rightline=false, leftline=true,
  leftmargin=10pt,
  skipabove=10pt, skipbelow=10pt,
  innertopmargin=5pt, innerbottommargin=5pt, innerleftmargin=10pt, innerrightmargin=10pt
}



% 创建包裹定理的环境
\surroundwithmdframed[style=theoremstyle]{theorem}
\surroundwithmdframed[style=theoremstyle]{lemma}
\surroundwithmdframed[style=theoremstyle]{corollary}
\surroundwithmdframed[style=theoremstyle]{proposition}

\surroundwithmdframed[style=definitionstyle]{definition}
\surroundwithmdframed[style=remarkstyle]{remark}
\surroundwithmdframed[style=exercisestyle]{exercise}
\surroundwithmdframed[style=exercisestyle]{problem}
